% !TeX program = pdflatex
\documentclass[12pt,a4paper]{article}

\newcommand{\notelanguage}{finnish}
% Shared preamble for course notes and exercise sets.

\usepackage[T1]{fontenc}
\usepackage{lmodern}
\usepackage[english]{babel}

\usepackage[a4paper,margin=30mm]{geometry}
\usepackage{microtype}
\usepackage{graphicx}
\usepackage{enumitem}

\usepackage{amsmath}
\usepackage{amssymb}
\usepackage{amsthm}
\usepackage{mathtools}
\usepackage{aliascnt}

\usepackage[hidelinks,pdfusetitle]{hyperref}

\numberwithin{equation}{section}
\setlist{itemsep=0.25em,topsep=0.5em}

\theoremstyle{plain}
\newtheorem{theorem}{Theorem}[section]
\newaliascnt{lemma}{theorem}
\newtheorem{lemma}[lemma]{Lemma}
\aliascntresetthe{lemma}
\newaliascnt{proposition}{theorem}
\newtheorem{proposition}[proposition]{Proposition}
\aliascntresetthe{proposition}
\newaliascnt{corollary}{theorem}
\newtheorem{corollary}[corollary]{Corollary}
\aliascntresetthe{corollary}

\theoremstyle{definition}
\newaliascnt{definition}{theorem}
\newtheorem{definition}[definition]{Definition}
\aliascntresetthe{definition}
\newaliascnt{example}{theorem}
\newtheorem{example}[example]{Example}
\aliascntresetthe{example}
\newaliascnt{problem}{theorem}
\newtheorem{problem}[problem]{Problem}
\aliascntresetthe{problem}

\theoremstyle{remark}
\newaliascnt{remark}{theorem}
\newtheorem{remark}[remark]{Remark}
\aliascntresetthe{remark}

\usepackage[nameinlink,noabbrev]{cleveref}

\crefname{theorem}{theorem}{theorems}
\crefname{lemma}{lemma}{lemmas}
\crefname{proposition}{proposition}{propositions}
\crefname{corollary}{corollary}{corollaries}
\crefname{definition}{definition}{definitions}
\crefname{example}{example}{examples}
\crefname{problem}{problem}{problems}
\crefname{remark}{remark}{remarks}

\DeclareMathOperator{\im}{im}
\DeclareMathOperator{\rank}{rank}
\DeclarePairedDelimiter{\abs}{\lvert}{\rvert}
\DeclarePairedDelimiter{\norm}{\lVert}{\rVert}
\DeclarePairedDelimiter{\set}{\lbrace}{\rbrace}

\usepackage{tikz}
\DeclareMathOperator{\sign}{sign}

\title{Johdatus yliopistomatematiikkaan\\Harjoituksen 3 ratkaisut}
\date{}
\author{}

\begin{document}

\exerciseheading{Johdatus yliopistomatematiikkaan}{Harjoituksen 3 ratkaisut}

\begin{exercise}
  Olkoon \(A = \set{-2, -1, 0, 1, 2}\). Ilmaise joukon \(A\) relaatio
  \(xRy \iff xy \geq 0\) luettelemalla sen alkiot eli kaikki siihen kuuluvat
  järjestetyt parit. Piirrä relaatiosta nuolikuvio. Määritä myös yhdistetty
  relaatio \(R^2 = R \circ R\).
\end{exercise}

\begin{solution}
  Määritellään merkkifunktio
  \[
    \sign(x) =
    \begin{cases}
      -1, & x < 0, \\
      0,  & x = 0, \\
      1,  & x > 0.
    \end{cases}
  \]
  Saadaan relaatiolle \(xRy\) ehto
  \[
    xRy \iff x = 0 \lor y = 0 \lor \sign(x) = \sign(y).
  \]
  Joten
  \[
    \begin{aligned}
      R = \bigl\{&
        (-2, -2), (-2, -1), (-2, 0), \\
        &(-1, -2), (-1, -1), (-1, 0), \\
        &(0, -2), (0, -1), (0, 0), (0, 1), (0, 2), \\
        &(1, 0), (1, 1), (1, 2), \\
        &(2, 0), (2, 1), (2, 2)
      \bigr\}.
    \end{aligned}
  \]
  Relaation nuolikuvio on
  \begin{center}
    \begin{tikzpicture}[
        >=stealth,
        element/.style={draw, circle, fill=white, minimum size=8mm,
        inner sep=0pt},
        relation/.style={->, thin}
      ]
      \draw (-2.8, 0) ellipse (1.0 and 3.0);
      \draw (2.8, 0) ellipse (1.0 and 3.0);
      \node at (-2.8, 3.35) {$A$};
      \node at (2.8, 3.35) {$A$};

      \node[element] (left-minus-two) at (-2.8, 2) {$-2$};
      \node[element] (left-minus-one) at (-2.8, 1) {$-1$};
      \node[element] (left-zero) at (-2.8, 0) {$0$};
      \node[element] (left-one) at (-2.8, -1) {$1$};
      \node[element] (left-two) at (-2.8, -2) {$2$};
      \node[element] (right-minus-two) at (2.8, 2) {$-2$};
      \node[element] (right-minus-one) at (2.8, 1) {$-1$};
      \node[element] (right-zero) at (2.8, 0) {$0$};
      \node[element] (right-one) at (2.8, -1) {$1$};
      \node[element] (right-two) at (2.8, -2) {$2$};

      \draw[relation] (left-minus-two) -- (right-minus-two);
      \draw[relation] (left-minus-two) -- (right-minus-one);
      \draw[relation] (left-minus-two) -- (right-zero);
      \draw[relation] (left-minus-one) -- (right-minus-two);
      \draw[relation] (left-minus-one) -- (right-minus-one);
      \draw[relation] (left-minus-one) -- (right-zero);
      \draw[relation] (left-zero) -- (right-minus-two);
      \draw[relation] (left-zero) -- (right-minus-one);
      \draw[relation] (left-zero) -- (right-zero);
      \draw[relation] (left-zero) -- (right-one);
      \draw[relation] (left-zero) -- (right-two);
      \draw[relation] (left-one) -- (right-zero);
      \draw[relation] (left-one) -- (right-one);
      \draw[relation] (left-one) -- (right-two);
      \draw[relation] (left-two) -- (right-zero);
      \draw[relation] (left-two) -- (right-one);
      \draw[relation] (left-two) -- (right-two);
    \end{tikzpicture}
  \end{center}
  Määritetään seuraavaksi relaatio \(R^2 = R \circ R\). Määritelmän mukaan
  \[
    R^2 = \set{(x, z) \in A^2 \colon \exists y \in A \colon (x, y) \in
    R \land (y, z) \in R}.
  \]
  Nyt, koska \(xR0\) ja \(0Rz\) kaikilla \(x,z \in A\), saadaan \(R^2
  = A \times A = A^2\).
\end{solution}

\begin{exercise}
  Jatkoa edelliseen tehtävään 1. Onko relaatio \(R\)
\begin{enumerate}[label=\alph*), nosep]
  \item refleksiivinen,
  \item irrefleksiivinen,
  \item symmetrinen,
  \item antisymmetrinen,
  \item asymmetrinen,
  \item transitiivinen,
  \item vertailullinen,
  \item ekvivalenssirelaatio?
\end{enumerate}
\end{exercise}

\begin{solution}
\leavevmode\par
\begin{enumerate}[label=\alph*), nosep]
\item On voimassa \(xRx\), sillä \(xx = x^2 \geq 0\) kaikilla \(x \in A\).
  \(R\) on siis \emph{refleksiivinen}.
\item \(R\) ei ole \emph{irrefleksiivinen}, sillä esimerkiksi \((0, 0) \in R\).
\item Näytetään \(xRy \implies yRx\) \emph{kaikilla} \(x,y \in A\).
  Koska \(xy = yx\) kaikilla \(x, y \in \ZZ\), saadaan
  \(xy \geq 0 \implies yx \geq 0\), eli \(R\) on \emph{symmetrinen}.
\item \((1, 0) \in R\) ja \((0, 1) \in R\), mutta \(0 \neq 1\). \(R\)
  ei siis ole \emph{antisymmetrinen}.
\item \((1, 0) \in R\) ja \((0, 1) \in R\), joten \(R\) ei ole
  \emph{asymmetrinen}.
\item \((-1, 0) \in R\) ja \((0, 1) \in R\), mutta \((-1, 1) \notin R\),
  sillä \((-1)1 = -1 < 0\). \(R\) ei siis ole \emph{transitiivinen}.
\item \((-1, 1) \notin R\) ja \((1, -1) \notin R\), joten \(R\) ei ole
  \emph{vertailullinen}.
\item \(R\) ei ole \emph{transitiivinen}, joten se ei ole
  \emph{ekvivalenssirelaatio}.
\end{enumerate}
\end{solution}

\begin{exercise}
Olkoot \(R\) ja \(S\) joukkojen \(A\) ja \(B\) välisiä relaatioita, toisin
sanoen \(R, S \subset A \times B\). Osoita, että
\begin{enumerate}[label=\alph*), nosep]
\item \((R^{-1})^{-1} = R\),
\item \((R \cap S)^{-1} = R^{-1} \cap S^{-1}\).
\end{enumerate}
\end{exercise}

\begin{solution}
\leavevmode\par
\begin{enumerate}[label=\alph*), nosep]
\item Olkoon \((a, b) \in A \times B\) mielivaltainen. Käänteisrelaation
määritelmän nojalla
\[
\begin{aligned}
  (a, b) \in (R^{-1})^{-1} &\iff (b, a) \in R^{-1} \\
  &\iff (a, b) \in R.
\end{aligned}
\]
Joukoilla \((R^{-1})^{-1}\) ja \(R\) on siis täsmälleen samat
alkiot, joten \((R^{-1})^{-1} = R\).

\item Olkoon \((a, b) \in B \times A\) mielivaltainen. Käänteisrelaation ja
joukon leikkauksen määritelmän nojalla
\[
\begin{aligned}
  (a, b) \in (R \cap S)^{-1} &\iff (b, a) \in (R \cap S) \\
  &\iff (b, a) \in R \land (b, a) \in S \\
  &\iff (a, b) \in R^{-1} \land (a, b) \in S^{-1} \\
  &\iff (a, b) \in R^{-1} \cap S^{-1}.
\end{aligned}
\]
Joukoilla \((R \cap S)^{-1}\) ja \(R^{-1} \cap S^{-1}\) on siis täsmälleen
samat alkiot, joten \((R \cap S)^{-1} = R^{-1} \cap S^{-1}\).
\end{enumerate}
\end{solution}

\begin{exercise}
Määritellään nollasta eroavien kokonaislukujen joukossa
\(\ZZ \setminus \set{0}\) relaatio
\[
x \sim y \iff \frac{x}{y} = \frac{y}{x}.
\]
Osoita, että \(\sim\) on ekvivalenssirelaatio.
\end{exercise}

\begin{solution}
Osoitetaan, että \(\sim\) on ekvivalenssirelaatio osoittamalla, että
\(\sim\) on refleksiivinen, symmetrinen ja transitiivinen.
\begin{enumerate}
\item Olkoon \(x \in \ZZ \setminus \set{0}\). Nyt
\(\frac{x}{x} = 1 = \frac{x}{x}\), joten \(x \sim x\). Relaatio \(\sim\) on
siis refleksiivinen.
\item Olkoot \(x,y \in \ZZ \setminus \set{0}\). Jos \(x \sim y\), niin
\(\frac{x}{y} = \frac{y}{x}\), joten myös
\(\frac{y}{x} = \frac{x}{y}\). Siis \(y \sim x\), joten relaatio on
symmetrinen.
\item Olkoot \(x,y,z \in \ZZ \setminus \set{0}\). Oletetaan, että
\(x \sim y\) ja \(y \sim z\). Koska \(x \sim y\), niin
\(\frac{x}{y} = \frac{y}{x} \implies x^2 = y^2\). Koska
\(y \sim z\), niin \(\frac{y}{z} = \frac{z}{y} \implies y^2 = z^2\).
Saadaan \(x^2 = z^2\), eli \(\frac{x}{z} = \frac{z}{x}\), joten
\(x \sim z\). Relaatio \(\sim\) on siis transitiivinen.
\end{enumerate}
On osoitettu, että \(\sim\) on refleksiivinen, symmetrinen ja transitiivinen.
Täten \(\sim\) on ekvivalenssirelaatio.
\end{solution}

\begin{exercise}
Jatkoa edelliseen tehtävään 4. Määritä kaikki ekvivalenssiluokat
\((\ZZ \setminus \set{0})/{\sim}\). Mikä on luvun \(3\) ekvivalenssiluokka
\([3]_{\sim}\)?
\end{exercise}

\begin{solution}
Olkoon \(x \in \ZZ \setminus \set{0}\). Luvun \(x\) ekvivalenssiluokka on
\[
[x]_{\sim} = \set{y \in \ZZ \setminus \set{0} \colon x \sim y}.
\]
Koska \(x \sim y \iff x^2 = y^2 \iff y = \pm x\), saadaan
\[
[x]_{\sim} = \set{-x, x}.
\]
Kaikki ekvivalenssiluokat ovat siis
\[
(\ZZ \setminus \set{0})/{\sim}
= \set{\set{-n, n} \colon n \in \ZZ_{+}},
\]
missä \(\ZZ_{+} = \set{1, 2, 3, \cdots}\) on positiivisten kokonaislukujen
joukko. Erityisesti luvun \(3\) ekvivalenssiluokka on
\[
[3]_{\sim} = \set{-3, 3}.
\]
\end{solution}

\begin{exercise}
Määritellään joukossa \([0, 4] \times \RR\) relaatio
\(xRy \iff y = \sqrt{x}\). Osoita, että \(R\) on kuvaus.

Huomaa, että merkintä \([a, b]\) tarkoittaa suljettua väliä
\(\set{x \in \RR : a \leq x \leq b}\), kun \(a, b \in \RR\), \(a < b\).
Käytetään nyt relaatiolle \(R\) tavallista funktiomerkintää
\(f : [0, 4] \to \RR\), \(f(x) = \sqrt{x}\).
\begin{enumerate}[label=\alph*), nosep]
\item Piirrä kuvauksen \(f\) kuvaaja.
\item Mitkä ovat kuvauksen \(f\) lähtö-, maali- ja arvojoukko?
\item Mikä on joukon \([1, 3]\) kuva? Entä mikä on joukon \([1, 3]\) alkukuva?
\item Onko kuvauksen \(f : [0, 4] \to \RR\) käänteisrelaatio \(f^{-1}\) kuvaus?
\end{enumerate}
\end{exercise}

\begin{solution}
Jokaisella \(x \in [0, 4]\) on \emph{täsmälleen yksi} epänegatiivinen
neliöjuuri \(\sqrt{x} \in \RR\). Siten jokaista lähtöjoukon alkiota vastaa
täsmälleen yksi maalijoukon alkio, joten \(R\) on kuvaus.
\leavevmode\par
\begin{enumerate}[label=\alph*), nosep]
\item Kuvauksen \(f(x) = \sqrt{x}\), \(x \in [0, 4]\), kuvaaja:
\begin{center}
\begin{tikzpicture}[x=1.7cm, y=1.7cm, >=stealth]
\draw[step=1, gray!25, very thin] (0, 0) grid (4, 2);
\draw[->] (-0.3, 0) -- (4.6, 0) node[right] {$x$};
\draw[->] (0, -0.3) -- (0, 2.7) node[above] {$y$};
\foreach \x in {1, 2, 3, 4}
\draw (\x, 0.04) -- (\x, -0.04) node[below] {$\x$};
\foreach \y in {1, 2}
\draw (0.04, \y) -- (-0.04, \y) node[left] {$\y$};
\node[below left] at (0, 0) {$0$};
\draw[thick, domain=0:2, samples=100, variable=\t]
plot ({\t*\t}, {\t});
\node[above] at (2.5, 1.65) {$y = \sqrt{x}$};
\fill (0, 0) circle (1.8pt);
\fill (4, 2) circle (1.8pt) node[above right] {$(4, 2)$};
\end{tikzpicture}
\end{center}
\item Lähtöjoukko on \([0, 4]\), maalijoukko on \(\RR\) ja arvojoukko
on \([0, 2]\).
\item Joukon \([1, 3]\) kuva on \(f([1, 3]) = [1, \sqrt{3}]\). Joukon
\([1, 3]\) alkukuva on
\[
\begin{aligned}
f^{-1}([1, 3]) &= \set{x \in [0, 4] \colon f(x) \in [1, 3]} \\
&= \set{x \in [0, 4] \colon 1 \leq \sqrt{x} \leq 3} \\
&= [0, 4] \cap [1, 9] = [1, 4].
\end{aligned}
\]
\item Funktion \(f\) käänteisrelaatio \(f^{-1} \subset \RR \times [0, 4]\)
ei ole kuvaus joukosta \(\RR\) joukkoon \([0, 4]\), sillä esimerkiksi
luvulle \(3\) ei ole kuvaa: millään \(x \in [0, 4]\) ei päde
\(\sqrt{x} = 3\).
Jos käänteisrelaation lähtöjoukoksi rajataan \([0, 2]\), saadaan kuvaus
\(g \colon [0, 2] \to [0, 4]\), \(g(y) = y^2\).
\end{enumerate}
\end{solution}

\end{document}
